% Compile with pdfLaTeX.
% Ensure member1.jpeg, member2.png, member3.jpeg, and architecture.png are in the same folder.

\documentclass[10pt,a4paper]{article}

% --- Page Geometry ---
\usepackage[
    a4paper,
    left=1.5cm,
    right=1.5cm,
    top=1.3cm,
    bottom=1.3cm,
    headheight=14pt,
    headsep=8pt,
    footskip=12pt
]{geometry}

% --- Fonts & Encoding (Times New Roman) ---
\usepackage[T1]{fontenc}
\usepackage[utf8]{inputenc}
\usepackage{newtxtext,newtxmath} % Complete Times New Roman font suite
\usepackage{microtype}

% --- Graphics & Colors ---
\usepackage{graphicx}
\usepackage{xcolor}
\definecolor{TitleBlue}{RGB}{41, 88, 142}
\definecolor{SubTextGray}{RGB}{85, 95, 110}
\definecolor{BorderGray}{RGB}{180, 190, 205}
\definecolor{BoxBg}{RGB}{252, 253, 255}
\definecolor{TitleBlack}{RGB}{0,0,0}

% --- Layout & Tables ---
\usepackage{array}
\usepackage{tabularx}
\usepackage{ragged2e}
\usepackage{enumitem}
\usepackage{fancyhdr}
\setlength{\headheight}{37.85225pt}
\usepackage{setspace}
\usepackage[most]{tcolorbox}
\usepackage[hidelinks,colorlinks=true,allcolors=TitleBlue]{hyperref}

% --- Paragraph Formatting ---
\setlength{\parindent}{0pt}
\setlength{\parskip}{0pt}
\renewcommand{\arraystretch}{1.15}

% --- Heading Commands ---
\newcommand{\bluebigheading}[1]{%
    {\bfseries\color{TitleBlue}\fontsize{13}{15}\selectfont #1}\par
}

\newcommand{\blueheading}[1]{%
    \vspace{2.5mm}%
    {\bfseries\color{TitleBlue}\fontsize{11}{13}\selectfont #1}\par
}

\newcommand{\instruction}[1]{%
    \vspace{0.5mm}%
    {\itshape\color{SubTextGray}\fontsize{8.5}{10.5}\selectfont #1}\par
    \vspace{1mm}%
}

% --- Styled Dynamic Answer Box ---
\newtcolorbox{dynamicanswerbox}[1][]{%
    enhanced,
    arc=1.5mm,
    colback=BoxBg,
    colframe=BorderGray,
    boxrule=0.6pt,
    left=3mm,
    right=3mm,
    top=2mm,
    bottom=2mm,
    boxsep=0mm,
    before skip=1.5mm,
    after skip=1.5mm,
    fontupper=\fontsize{8.5}{11}\selectfont\normalfont,
    #1
}

% --- Header & Footer ---
\pagestyle{fancy}
\fancyhf{}
\renewcommand{\headrulewidth}{0.4pt}
\renewcommand{\headrule}{\hbox to\headwidth{\color{BorderGray}\leaders\hrule height \headrulewidth\hfill}}
\renewcommand{\footrulewidth}{0pt}

\fancyhead[L]{%
    \bfseries\fontsize{8.5}{10}\selectfont\color{TitleBlack}
    \textbf{DS with C (TCS-302) \& OOPS with C++ (TCS-307) \\ DSCPP-III-2026-T224}
    
    
}
\fancyhead[R]{%
    \bfseries\fontsize{8.5}{10}\selectfont\color{TitleBlue}
    PROJECT PROPOSAL
}
\fancyfoot[R]{%
    \fontsize{8}{9.5}\selectfont\color{SubTextGray}
    Page \thepage
}

\begin{document}

% =================================================
% PAGE 1
% =================================================
\bluebigheading{PROJECT AND TEAM INFORMATION}
\vspace{2mm}
\textbf{VeriNews is the model we propose to solve the problem.\\}
\blueheading{\\Project Title}

\begin{dynamicanswerbox}
\textbf{VeriNews: An Explainable Multi-Factor News Verification and Credibility Scoring Engine Using DSA in C/C++}
\end{dynamicanswerbox}
\blueheading{Student / Team Information}
\vspace{1.5mm}

% Team information table
\noindent
\begin{tabularx}{\textwidth}{|>{\RaggedRight\arraybackslash}p{0.52\textwidth}|>{\RaggedRight\arraybackslash}X|}
\hline
\textbf{Team Name:} ByteForce & \textbf{Team \#:} 07 \\
\hline
\begin{minipage}[c][3.8cm][c]{\linewidth}
    \textbf{Team member 1 (Team Lead)}\\[1mm]
    \textit{Srivastava, Shambhavi}\\
    \textbf{ID:} 2510012471\\
    \textbf{Email:} \href{mailto:2510012471@geu.ac.in}{2510012471@geu.ac.in}
\end{minipage}
&
\begin{minipage}[c][3.8cm][c]{\linewidth}
    \centering
    \includegraphics[width=2.5cm,height=3.2cm,keepaspectratio]{member1.jpeg}
\end{minipage}
\\
\hline
\begin{minipage}[c][3.8cm][c]{\linewidth}
    \textbf{Team member 2}\\[1mm]
    \textit{Panwar, Abhinav}\\
    \textbf{ID:} 2512510050\\
    \textbf{Email:} \href{mailto:2512510050@geu.ac.in}{2512510050@geu.ac.in}
\end{minipage}
&
\begin{minipage}[c][3.8cm][c]{\linewidth}
    \centering
    \includegraphics[width=2.5cm,height=3.2cm,keepaspectratio]{member2.png}
\end{minipage}
\\
\hline
\begin{minipage}[c][3.8cm][c]{\linewidth}
    \textbf{Team member 3}\\[1mm]
    \textit{Sharma, Akshat Mohan}\\
    \textbf{ID:} 2510010700\\
    \textbf{Email:} \href{mailto:2510010700@geu.ac.in}{2510010700@geu.ac.in}
\end{minipage}
&
\begin{minipage}[c][3.8cm][c]{\linewidth}
    \centering
    \includegraphics[width=2.5cm,height=3.2cm,keepaspectratio]{member3.jpeg}
\end{minipage}
\\
\hline
\end{tabularx}

\newpage

% =================================================
% PAGE 2
% =================================================
\bluebigheading{PROPOSAL DESCRIPTION (10 pts)}
\vspace{5mm}
\blueheading{Motivation (1 pt)}

\begin{dynamicanswerbox}
In today’s digital environment, false news and exaggerated headlines spread quickly on social media and messaging apps.
These misleading stories can change people's beliefs, create fear, and damage trust in real news.It is difficult for individuals and students to verify the truth of a headline or article as it appears due to the fast pace and large amount of unfiltered information.\\
Most existing fact-checking tools either depend on manual checks that are slow or use complex machine learning models that are not transparent.
The real challenge is not just labeling news as fake or real, but providing users with a credibility score that gives a clear explanation of why a news item was flagged.\\To address this, we are building \textbf{VeriNews}, a fast and explainable news verification engine using fundamental data structures in C and object-oriented programming in C++.
By identifying key entities, matching them with verified or questionable records from curated databases, and calculating a score based on source reliability, topic relevance, and time consistency, our platform offers an easy-to-use, educational, and computationally efficient solution to help tackle misinformation from its source.
\end{dynamicanswerbox}
\vspace{10mm}
\blueheading{State of the Art / Current solution (1 pt)}

\begin{dynamicanswerbox}
Currently, fake news verification is mostly limited to two main approaches.
One is manual fact-checking by experts, as seen in platforms like Snopes, Alt News, and PolitiFact.These are slow but accurate.On the other hand, automated systems use algorithms to detect false news quickly.However, they often lack the ability to understand context and nuances.As a result, users are forced to choose between slow but accurate methods or fast but unreliable ones.There is very little in between.Many systems also rely on cloud-based AI or NLP APIs, which are not transparent.Manual methods take too long, allowing fake news to spread before it can be corrected.
On the technology front, most platforms use complex neural networks or proprietary cloud-based fact-checking tools.\\
These systems work like black boxes, giving users no insight into how they make decisions or how they evaluate sources.Also, existing systems rarely provide simple, standalone rule-based software that can run locally with efficient algorithms, without needing many external dependencies.
\end{dynamicanswerbox}
\vspace{10mm}
\blueheading{Project Goals and Milestones (2 pts)}

\begin{dynamicanswerbox}
Our goal is to build a high-performance verification and credibility-scoring application that allows users to submit news headlines or article snippets and immediately receive a weighted credibility score (0–100\%) along with a detailed explanation.

\vspace{1mm}
\textbf{Project Milestones:}
\begin{itemize}[leftmargin=4.5mm,itemsep=1.5pt,topsep=1.5pt,partopsep=0pt,parsep=0pt]
    \item \textbf{Milestone 1 (Phase 1):}  Implement custom C data structures: Hash Tables (indexing), Linked Lists (records), and BSTs (sorting/filtering).
    \item \textbf{Milestone 2 (Phase 2):}  Build string tokenization, stop-word filtering, entity extraction, and the weightedcredibility scoring engine.
    \item \textbf{Milestone 3 (Phase 2/3):}  Implement C++ OOP architecture (User, News,Admin, Engine),a FIFOrequest queue, and anundo/history stack.
    \item \textbf{Milestone 4 (Phase 3):}  Build file I/O for datapersistence and theinteractiveAdmin CRUD module.
    \item \textbf{Milestone 5 (Phase 3):}  Conduct unit testing, complexity profiling, and compile final submission deliverables.
\end{itemize}
\end{dynamicanswerbox}
\vspace{10mm}
\blueheading{Project Approach (3 pts)}

\begin{dynamicanswerbox}
The system combines performance-oriented data structures with an object-oriented C++ control layer.

\vspace{1mm}
\begin{itemize}[leftmargin=4.5mm,itemsep=1.5pt,topsep=1.5pt,partopsep=0pt,parsep=0pt]
    \item \textbf{Text Processing \& Keyword Extraction:}  Cleaninput text, removeunnecessaryelements, filter stop words, and extract key entities.
    \item \textbf{Data Structures \& Matching Layer (C):} HashTablesareusedforquickkeywordlookupsinnearconstanttime,dynamicLinkedListsareused
to chain records, Binary Search Trees (BST) are used for filtering based on time or score ranges, FIFO Queues are used for grouping input data,
and Stacksareused to track operation history.
    \item \textbf{Multi-Factor Verification Engine:} This system calculates a score based on several factors: Entity Match (40\%) Source Reliability (25\%) Corroborating Reports (20\%) Temporal Consistency (10\%) and Category Alignment (5\%).
    \item \textbf{Classification \& Persistence (C++ OOP):} Output is categorized into three groups: Likely Verified (80–100\%) Needs Verification (\%) or Suspicious (0–49\%).
\end{itemize}
\end{dynamicanswerbox}

\newpage

% =================================================
% PAGE 3
% =================================================
\vspace{7mm}
\blueheading{System Architecture (High Level Diagram) (2 pts)}

\begin{dynamicanswerbox}
\centering
\vspace{1mm}
\includegraphics[width=0.96\textwidth,height=15cm,keepaspectratio]{architecture.png}
\vspace{1mm}
\end{dynamicanswerbox}
\vspace{7mm}
\blueheading{Project Outcome / Deliverables (1 pt)}

\begin{dynamicanswerbox}
By the end of this project, we will deliver a fully working, standalone C/C++ application called \textbf{VeriNews} that allows users to assess news headlines and view detailed credibility scores.\\ Key Deliverables include:
\begin{itemize}[leftmargin=4.5mm,itemsep=1pt,topsep=1pt,partopsep=0pt,parsep=0pt]
    \item Source Code: Fully modular C (custom ADTs) and C++ (classes and control flow) code.
    \item Dataset \& Persistence Layer: Seeded news corpus, trusted source registry, and file-based local storage.
    \item Admin Module: Interactive console interface for full database CRUD operations.
    \item Documentation: Final project report with Big-O complexity analysis and UML diagrams.
\end{itemize}
\end{dynamicanswerbox}
\vspace{7mm}
\blueheading{Assumptions}

\begin{dynamicanswerbox}
Input text is English and provides sufficient named entities for tokenization.
Evaluations rely on a curated local truth dataset; it is an offline verification engine, not a live web crawler.
\end{dynamicanswerbox}
\vspace{7mm}
\blueheading{References}

\begin{dynamicanswerbox}
\begin{enumerate}[leftmargin=5mm,label={[\arabic*]},itemsep=0.5pt,topsep=0.5pt,partopsep=0pt,parsep=0pt]
    \item Vosoughi, S., Roy, D., \& Aral, S. (2018). The spread of false news online. \textit{Science}, 359(6380), 1146--1151. \url{https://doi.org/10.1126/science.aap9559}
    \item Shu, K., Sliva, A., Wang, S., Tang, J., \& Liu, H. (2017). Fake News Detection on Social Media. \textit{ACM SIGKDD Explorations}, 19(1), 22--36. \url{https://doi.org/10.1145/3137597.3137600}
    \item Cormen, T. H., et al. (2009). \textit{Introduction, to Algorithms} Third edition. MIT Press.
    \item Stroustrup, B. (2013). \textit{The C++ Programming Language} (Fourth edition. Addison-Wesley.
    \item Lazer, D. M. And others. (2018). The science of news.  \textit{Science}, 359(6380), 1094-1096. \url{https://doi.org/10.1126/science.aao2998}
\end{enumerate}
\end{dynamicanswerbox}

\end{document}
